% SGA 3, Expose III, Loop 1 English reconstruction.
% Scope: section 2, from the heading through Corollary 2.14, stopping
% immediately before section 3.
% Authority: Polo--Gille Exp3-14oct24.pdf, component-local pp. 31--37
% (combined-reader pp. 141--147; printed pp. 131--137), with only the
% section-2 portions of the first and last pages admitted. Force-OCR is a
% locator only. English wording was compared with Jacob C. Reinhold's CC BY
% 4.0 Markdown lineage at commit e7a259f3f8608ad3edf9bf6eead3fd504dd2d23e.
% The Polo--Gille PDF controls prose, mathematics, numbering, and notes.

\setcounter{footnote}{51}

% Source page: Exp3 local p. 31; combined p. 141; printed p. 131.
\section*{\SGAInlineRefTarget{sga3:iii:section:2}2. 群概形态射的无穷小扩张}
\addcontentsline{toc}{section}{2. 群概形态射的无穷小扩张}

沿用 \S0 的记号. 设 \(Y\) 和 \(X\) 是两个 \(S\)-群函子,
并设 \(M\) 是群同态 \(p_X:X\to X^+\) 的核.
于是有 \(S\)-群函子的正合列
\[
 1\longrightarrow M\longrightarrow X\xrightarrow{p_X}X^+.
\]
由 \(X^+\) 的定义, 有同构
\begin{align*}
 \operatorname{Hom}_S(Y,X^+)
  &\xrightarrow{\sim}\operatorname{Hom}_{S_{\mathscr J}}(Y_{\mathscr J},X_{\mathscr J}),\\
 \operatorname{Hom}_{S\text{-}\mathrm{gr}}(Y,X^+)
  &\xrightarrow{\sim}
    \operatorname{Hom}_{S_{\mathscr J}\text{-}\mathrm{gr}}
    (Y_{\mathscr J},X_{\mathscr J}).
\end{align*}
映射
\[
 \operatorname{Hom}_S(Y,p_X):
 \operatorname{Hom}_S(Y,X)\longrightarrow\operatorname{Hom}_S(Y,X^+)
\]
将 \(S\)-态射 \(f:Y\to X\) 映到 \(S\)-态射
\(f^+:Y\to X^+\), 后者在这些同构下对应于由 \(f\) 作基变换得到的
\(S_{\mathscr J}\)-态射
\(f_{\mathscr J}:Y_{\mathscr J}\to X_{\mathscr J}\).
若 \(M\) 交换, 则 \SGARefLink{sga3:iii:proposition:1-2-4}{命题 1.2.4}
适用于这一情形.

\medskip
\noindent\SGAInlineRefTarget{sga3:iii:paragraph:2-0}\textbf{2.0.}
\footnote{编者注: 为便于后文引用, 新增了编号
\SGARefLink{sga3:iii:paragraph:2-0}{2.0}.}
以下考虑 \(Y\) 在 \(S\) 上平坦, 而 \(X\) 为下列两类之一的
\(S\)-群函子的情形:
\begin{enumerate}
\item[\(a\)] \(X\) 是 \(S\)-群概形;
\item[\(b\)]\SGAInlineRefTarget{sga3:iii:paragraph:2-0:item:b}
  \(X=\underline{\operatorname{Aut}}_S(E)\), 其中 \(E\) 是
  \(S\) 上有限呈示的 \(S\)-概形.
\end{enumerate}
用 \((\mathrm{Flat})/S\) 表示在 \(S\) 上平坦的
\(S\)-概形所成的范畴. 在情形 \(a\), 分别地, 情形 \(b\) 中,
\(S\)-群函子 \(M=\operatorname{Ker}(X\to X^+)\),
它在 \((\mathrm{Flat})/S\) 上的限制 \(L\),
以及 \(X\) 的内自同构在 \(M\) 上的作用,
已在 \SGARefLink{sga3:iii:corollary:0-9}{0.9},
\SGARefLink{sga3:iii:remark:0-5}{0.5} 和 \SGARefLink{sga3:iii:corollary:0-10}{0.10},
分别地, 在 \SGARefLink{sga3:iii:proposition:0-11}{0.11},
\SGARefLink{sga3:iii:remark:0-14}{0.14} 和 \SGARefLink{sga3:iii:corollary:0-12}{0.12}
中算出.

更明确地说, 在情形 \(a\) 中, 设 \(L_0\) 是取值于交换群的
\(S_0\)-函子, 它对每个 \(S_0\)-概形 \(T_0\) 由下式定义:
\[
 \operatorname{Hom}_{S_0}(T_0,L_0)
 =\operatorname{Hom}_{\mathscr O_{T_0}}
 \left(
  \omega^1_{X_0/S_0}\otimes_{\mathscr O_{S_0}}\mathscr O_{T_0},
  \mathscr J\otimes_{\mathscr O_{S_0}}\mathscr O_{T_0}
 \right).
\]
群 \(X_0\) 通过它在 \(\omega^1_{X_0/S_0}\) 上的伴随表示
作用于此函子, 并且对每个 \(S\)-概形 \(T\) 有
\[
 L=\prod_{S_0/S}L_0;
 \qquad L(T)=L_0(T\times_S S_0)
\]

在情形 \(b\) 中, 设 \(\pi:E\to S\) 是结构态射.
则 \(L\) 是 \((\mathrm{Flat})/S\) 上的交换群函子, 由
\[
 \operatorname{Hom}_S(T,L)
 =\Gamma\!\left(
  T,
  \pi_*\operatorname{Hom}_{\mathscr O_E}
   (\Omega^1_{E/S},\mathscr J\mathscr O_E)
  \otimes_{\mathscr O_S}\mathscr O_T
 \right),
\]
定义, 而将 \(X\) 看成 \((\mathrm{Flat})/S\) 上的函子时,
它如 \SGARefLink{sga3:iii:corollary:0-12}{0.12} 中所述作用于它.
在两种情形中, 都有 \((\mathrm{Flat})/S\) 上的群函子正合列:
\[
 \tag{E}\label{sga3:iii:paragraph:2-0:equation:e}
 1\longrightarrow L\longrightarrow X\longrightarrow X^+.
\]

由于 \(Y\) 在 \(S\) 上平坦, \SGARefLink{sga3:iii:lemma:1-1-1}{引理 1.1.1}
表明, 群 \(H^i(Y,M)\) 仅依赖于 \(M\) 在
\((\mathrm{Flat})/S\) 上的限制 \(L\). 在情形 \(a\) 中,
\(L=\prod_{S_0/S}L_0\), \SGARefLink{sga3:iii:lemma:1-1-2}{引理 1.1.2}
给出同构
\[
 H^i(Y,L)\simeq H^i(Y_0,L_0).
\]

% Source pages: Exp3 local pp. 32--33; combined pp. 142--143;
% printed pp. 132--133.
综合以上各点, \SGARefLink{sga3:iii:proposition:1-2-4}{命题 1.2.4}
给出如下结果.
\footnote{编者注: 原来放在 \SGARefLink{sga3:iii:theorem:2-1}{定理 2.1}
陈述中并占据原版第 113 页大部分篇幅的定义, 已移到前面.}

\medskip
\noindent\SGAInlineRefTarget{sga3:iii:theorem:2-1}\textbf{定理 2.1.}
\textit{设 \(S\) 是概形, \(\mathscr I\supset\mathscr J\)
是两个满足 \(\mathscr I\mathscr J=0\) 的拟凝聚理想,
分别定义闭子概形 \(S_0\) 和 \(S_{\mathscr J}\). 设}
\begin{itemize}
\item[] \textit{\(X\) 是类型 \(a\) 或 \(b\) 的 \(S\)-群函子,
  而 \(L_0\) 和 \(L\) 如上所述;}
\item[] \textit{\(Y\) 是在 \(S\) 上平坦的 \(S\)-群概形,
  \(f_{\mathscr J}:Y_{\mathscr J}\to X_{\mathscr J}\) 是
  \(S_{\mathscr J}\)-群态射.}
\end{itemize}
\textit{则:}
\begin{enumerate}
\item[\(i\)] \textit{态射 \(f_{\mathscr J}\) 可提升为
  \(S\)-群态射 \(Y\to X\), 当且仅当下列两个条件成立:}
  \begin{enumerate}
  \item[\(i_1\)] \textit{\(f_{\mathscr J}\) 可提升为
    \(S\)-函子态射 \(Y\to X\). 由 \SGARefLink{sga3:iii:proposition:1-2-4}{命题 1.2.4},
    这定义了 \(Y\) 在 \(L\) 上的作用, 且与所选提升无关.
    在情形 \(a\) 中, 所得的 \(Y_0\) 在 \(L_0\) 上的作用来自
    \(f_0:Y_0\to X_0\) 及 \(X_0\) 在 \(L_0\) 上的伴随作用.}
  \item[\(i_2\)] \textit{典范定义的障碍
    \(c(f_{\mathscr J})\in H^2(Y,L)\) 为零.
    在情形 \(a\) 中, 此群同构于 \(H^2(Y_0,L_0)\).}
  \end{enumerate}
\item[\SGAInlineRefTarget{sga3:iii:theorem:2-1:item:ii}(ii)]
  \textit{若 \(i\) 中的条件成立, 则延拓 \(f_{\mathscr J}\) 的
  \(S\)-群函子态射 \(Y\to X\) 所成的集合 \(\mathcal E\)
  在 \(Z^1(Y,L)\) 下主齐性, 后者在情形 \(a\) 中同构于
  \(Z^1(Y_0,L_0)\). 此外, 将 \(\mathcal E\) 模去
  \(X\) 的由如下截面定义的内自同构的作用: 这些 \(S\) 上的截面
  在 \(S_{\mathscr J}\) 上诱导 \(X_{\mathscr J}\) 的单位截面;
  所得商集在 \(H^1(Y,L)\) 下主齐性, 后者在情形 \(a\)
  中同构于 \(H^1(Y_0,L_0)\).}
\item[\SGAInlineRefTarget{sga3:iii:theorem:2-1:item:iii}(iii)]
  \textit{若 \(f:Y\to X\) 是延拓 \(f_{\mathscr J}\) 的
  \(S\)-群函子态射, 则通过这些内自同构由 \(f\)
  得到的态射 \(Y\to X\) 所成的集合同构于}
  \[
   \Gamma(L)/H^0(Y,L),
  \]
  \textit{并且在情形 \(a\) 中同构于
  \(\Gamma(L_0)/H^0(Y_0,L_0)\).}
\end{enumerate}

\medskip
\noindent\SGAInlineRefTarget{sga3:iii:remark:2-1-1}\textbf{注 2.1.1.}
\footnote{编者注: 此注与 \SGARefLink{sga3:iii:remark:4-5-1}{4.5.1} 类似,
是为引入 \SGARefLink{sga3:iii:paragraph:4-38}{4.38} 中所用的记号
\(d(f,f')\) 及 \(\bar d(f,f')\) 而新增的.
因此, \SGARefLink{sga3:iii:theorem:2-1:item:ii}{2.1(ii)}
中关于 \(\mathcal E\) 本身的部分也是新增的.}
若 \(f,f':Y\to X\) 是延拓 \(f_{\mathscr J}\) 的
\(S\)-群函子态射, 则得到一个余循环
\[
 d(f,f')\in Z^1(Y,L)
 \quad\bigl(\simeq Z^1(Y_0,L_0)\text{ 在情形 (a) 中}\bigr)
\]
使得
\begin{equation}
 f'=d(f,f')\cdot f.
 \tag{*}\label{sga3:iii:remark:2-1-1:equation:star}
\end{equation}
\footnote{编者注: 我们遵循原文的符号约定. 也可以写成
\(f'=d(f',f)\cdot f\), 但如此一来, 为对两个浸入
\(i,i':Y\hookrightarrow X\) 保留 \SGARefLink{sga3:iii:lemma:4-27}{4.27}
中的等式 \(d_1(d(i,i'))=d(Y,i'(Y))\),
就必须改变 \SGARefLink{sga3:iii:remark:4-5-1}{4.5.1}
中引入的类 \(d(Y,Y')\) 的符号. 这又会改变
\SGARefLink{sga3:iii:proposition:4-8}{4.8},
\SGARefLink{sga3:iii:corollary:4-14}{4.14} 和 \SGARefLink{sga3:iii:lemma:4-17}{4.17}
中的符号. 我们选择保留原文的全部符号, 它们是正确的.}
用 \(\bar d(f,f')\) 表示 \(d(f,f')\) 在
\(H^1(Y,L)\) 中的像, 在情形 \(a\) 中等价地表示它在
\(H^1(Y_0,L_0)\) 中的像.

% Source page: Exp3 local p. 33; combined p. 143; printed p. 133.
\medskip
\noindent\SGAInlineRefTarget{sga3:iii:remark:2-2}\textbf{注 2.2.}
沿用前述记号; 特别地, \(Y\) 在 \(S\) 上平坦.
在情形 \(b\) 中, 公式 (\SGARefLink{sga3:iii:remark:0-14:equation:0-14-1}{0.14.1})
将 \(L\) 与
\[
 \mathbf W\!\left(
  \pi_*\operatorname{Hom}_{\mathscr O_E}
  (\Omega^1_{E/S},\mathscr J\mathscr O_E)
 \right).
\]
在 \((\mathrm{Flat})/S\) 上的限制等同.
\SGAInlineRefTarget{sga3:iii:remark:2-2:item:a}
在情形 \(a\) 中, 再假设 \(X\) 在 \(S\) 上局部有限呈示.
则公式 (\SGARefLink{sga3:iii:remark:0-6:equation:0-6-1}{0.6.1})
将 \(L\) 与
\[
 \prod_{S_0/S}\mathbf W\!\left(
  \operatorname{Hom}_{\mathscr O_{S_0}}
   (\omega^1_{X_0/S_0},\mathscr J)
 \right).
\]
在 \((\mathrm{Flat})/S\) 上的限制等同.
在两种情形中, 由 EGA I, 9.1.1, 用来构造 \(\mathbf W\) 的模都是拟凝聚的.
再假设 \(Y\) 在 \(S\) 上仿射.
\footnote{编者注: 此假设及对 I, \SGARefLink{sga3:i:subsection:5-3}{5.3}
的引用是新增的.} 由 I, \SGARefLink{sga3:i:subsection:5-3}{5.3}, 得到
\begin{align*}
 H^i(Y,L)=H^i(Y_0,L_0)
  &=H^i\!\left(Y_0,
    \operatorname{Hom}_{\mathscr O_{S_0}}
    (\omega^1_{X_0/S_0},\mathscr J)\right)
  &&\text{在情形 (a) 中},\\
 H^i(Y,L)
  &=H^i\!\left(Y,
    \pi_*\operatorname{Hom}_{\mathscr O_E}
    (\Omega^1_{E/S},\mathscr J\mathscr O_E)\right)
  &&\text{在情形 (b) 中}.
\end{align*}

% Source page: Exp3 local p. 34; combined p. 144; printed p. 134.
\medskip
\noindent\SGAInlineRefTarget{sga3:iii:remark:2-3}\textbf{注 2.3.}
\begin{enumerate}
\item[(1)] 由 \SGARefLink{sga3:iii:proposition:0-16}{0.16} 和
  \SGARefLink{sga3:iii:corollary:0-17}{0.17}, 当 \(Y\) 仿射且
  \[
   \begin{cases}
    X\text{ 在 }S\text{ 上平滑}, &\text{在情形 (a) 中},\\
    E\text{ 在 }S\text{ 上平滑且仿射}, &\text{在情形 (b) 中}.
   \end{cases}
  \]
  时, 条件 \(i_1\) 自动成立.
\item[(2)] 在这些假设下, \(Y\) 仍如
  \SGARefLink{sga3:iii:paragraph:2-0}{2.0} 中那样在 \(S\) 上平坦,
  则 \SGARefLink{sga3:iii:remark:2-2:item:a}{公式 2.2(a)} 和
  (\SGARefLink{sga3:iii:remark:0-6:equation:0-6-2}{0.6.2}) 在情形 \(a\) 中给出
  \[
   H^i(Y,L)=H^i(Y_0,L_0)
   =H^i\!\left(Y_0,
    \operatorname{Lie}(X_0/S_0)
    \otimes_{\mathscr O_{S_0}}\mathscr J\right),
  \]
  而在情形 \(b\) 中, 由 (\SGARefLink{sga3:iii:remark:0-14:equation:0-14-2}{0.14.2}),
  \SGARefLink{sga3:iii:lemma:1-1-2}{1.1.2} 和 I, \SGARefLink{sga3:i:subsection:5-3}{5.3},
  \[
   H^i(Y,L)=H^i\!\left(Y_0,
    \pi_{0*}\operatorname{Hom}_{\mathscr O_{E_0}}
    (\Omega^1_{E_0/S_0},
     \mathscr J\otimes_{\mathscr O_{S_0}}\mathscr O_{E_0})\right).
  \]
  \footnote{编者注: 前面的内容是新增的; 见
  \SGARefLink{sga3:iii:remark:0-14}{0.14} 中的编者注 31.}
\end{enumerate}

现在对可对角化 \(S\)-群 \(Y\) (I, \SGARefLink{sga3:i:subsection:4-4}{4.4})
陈述几个推论. 若 \(S\) 仿射, 则由 I, \SGARefLink{sga3:i:theorem:5-3-3}{5.3.3},
对每个 \(n>0\) 及每个拟凝聚 \(\mathscr O_S\)-模 \(F\),
都有 \(H^n(Y,F)=0\).

\medskip
\noindent\SGAInlineRefTarget{sga3:iii:corollary:2-4}\textbf{推论 2.4.}
\textit{设 \(S_0\hookrightarrow S\) 是由幂零理想定义的闭子概形,
\(Y\) 是可对角化 \(S\)-群. 设}
\begin{enumerate}
\item[\(a\)] \textit{\(X\) 是在 \(S\) 上局部有限呈示的
  \(S\)-群, 或者}
\item[\(b\)]\SGAInlineRefTarget{sga3:iii:corollary:2-4:item:b}
  \textit{\(X=\underline{\operatorname{Aut}}_S(E)\), 其中
  \(E\) 是局部有限呈示的 \(S\)-概形.}
\end{enumerate}
\textit{若 \(f:Y\to X\) 是 \(S\)-群态射,
其基变换 \(f_0:Y_0\to X_0\) 是单位态射, 则 \(f\) 是单位态射.}

事实上, 问题在 \(S\) 上是局部的, 在情形 \(b\) 中也在 \(E\) 上是局部的.
因此可假设 \(S\) 仿射, 并在情形 \(b\) 中假设
\(E\) 在 \(S\) 上有限呈示. 引入闭子概形 \(S_n\subset S\),
它们由定义 \(S_0\) 的理想的各次幂定义.
这将断言归结为 \(S_0\) 由平方为零的理想定义的情形,
此时断言由 \SGARefLink{sga3:iii:theorem:2-1}{定理 2.1} 和
\SGARefLink{sga3:iii:remark:2-2}{注 2.2} 得到.

若不要求 \(f_0\) 为单位态射, 则有如下陈述.

\medskip
\noindent\SGAInlineRefTarget{sga3:iii:corollary:2-5}\textbf{推论 2.5.}
\textit{设 \(S\) 和 \(S_0\) 如 \SGARefLink{sga3:iii:corollary:2-4}{2.4}
中所述, 并且 \(S\) 仿射. 设 \(Y\) 是可对角化 \(S\)-群,
\(X\) 是 \(S\)-群函子, \(f_0:Y_0\to X_0\) 是
\(S_0\)-群函子态射.}
\begin{enumerate}
\item[\(i\)]\footnote{编者注: 原文陈述为:
``假设下列两个性质之一成立: \(a\) \(X\) 是局部有限呈示的
\(S\)-群; \(b\)
\(X=\underline{\operatorname{Aut}}_S(E)\), 其中 \(E\) 在 \(S\)
上有限呈示. 则 \(f_0\) 可延拓为 \(S\)-群态射
\(Y\to X\), 当且仅当它可延拓为 \(S\)-函子态射 \(Y\to X\).''
原文还说, 这由定理的 \SGARefLink{sga3:iii:corollary:2-5:item:ii}{(ii)}
项逐步推出. 这个论证似乎不充分. 例如, 若
\(\mathscr I^3=0\) 且 \(\mathscr J=\mathscr I^2\),
并且 \(f:Y\to X\) 是提升 \(f_0\) 的 \(S\)-函子态射,
则 \(f_0\) 可提升为一个 \(S_{\mathscr J}\)-群态射
\(g_{\mathscr J}:Y_{\mathscr J}\to X_{\mathscr J}\);
但 \(g_{\mathscr J}\) 是否必可再提升为 \(S\)-函子态射 \(g:Y\to X\)?
无论如何, 断言 \(i\) 在后文并不使用, 后文始终假定 \(X\) 在 \(S\) 上平滑.}
\item[\SGAInlineRefTarget{sga3:iii:corollary:2-5:item:ii}(ii)]
\textit{假定 \(X\) 是平滑 \(S\)-群, 或者
\(X=\underline{\operatorname{Aut}}_S(E)\), 且 \(E\) 在 \(S\) 上平滑且仿射.
则 \(f_0\) 可延拓为 \(S\)-群态射 \(Y\to X\),
且任意两个延拓都在 \(X\) 的某个内自同构下共轭,
此内自同构由一个 \(S\) 上的截面定义,
该截面在 \(S_0\) 上诱导 \(X_0\) 的单位截面.}
\end{enumerate}

% Source page: Exp3 local p. 35; combined p. 145; printed p. 135.
如上引入 \(S_n\).
\footnote{编者注: 下面略微修改并扩充了原文.}
对于 \SGARefLink{sga3:iii:corollary:2-5:item:ii}{(ii)},
在 \(S\) 上平滑的概形必然局部有限呈示.
因此, 在情形 \(b\) 中, 平滑且仿射的 \(E/S\) 为有限呈示,
所以 \SGARefLink{sga3:iii:paragraph:2-0:item:b}{2.0(b)} 的假设确实满足.

由 \SGARefLink{sga3:iii:proposition:0-16}{0.16} 和
\SGARefLink{sga3:iii:corollary:0-17}{0.17},
\SGARefLink{sga3:iii:theorem:2-1}{定理 2.1} 的条件 \(i_1\) 自动成立.
此外, \(X_{S_n}\) 在 \(S_n\) 上的每个截面均可提升为
\(X_{S_{n+1}}\) 在 \(S_{n+1}\) 上的截面:
在情形 \(a\) 中这由平滑性得到,
在情形 \(b\) 中由 \SGARefLink{sga3:iii:corollary:0-17}{0.17} 得到.
因此, 若 \(f\) 和 \(f'\) 是 \(f_0\) 的两个提升,
则通过提升 \SGARefLink{sga3:iii:theorem:2-1:item:ii}{定理 2.1(ii)}
所给出的内自同构, 可逐步使 \(f_n=f'_n\). 这就证明了断言.

同一论证结合 \SGARefLink{sga3:iii:remark:2-3}{注 2.3}, 给出:

\medskip
\noindent\SGAInlineRefTarget{sga3:iii:corollary:2-6}\textbf{推论 2.6.}
\textit{设 \(\mathscr I\) 是定义 \(S_0\) 的 \(S\) 的幂零理想,
\(Y\) 是在 \(S\) 上平坦的仿射 \(S\)-群,
\(X\) 是平滑 \(S\)-群.}
\begin{enumerate}
\item[\(i\)] \textit{若对每个 \(n\geq0\), 有}
 \[
  H^2\!\left(Y_0,
   \operatorname{Lie}(X_0/S_0)\otimes_{\mathscr O_{S_0}}
   \mathscr I^{n+1}/\mathscr I^{n+2}\right)=0,
 \]
 \textit{则每个 \(S_0\)-群态射 \(f_0:Y_0\to X_0\)
 都可提升为 \(S\)-群态射 \(f:Y\to X\).}
\item[(ii)] \textit{若对每个 \(n\geq0\), 有}
 \[
  H^1\!\left(Y_0,
   \operatorname{Lie}(X_0/S_0)\otimes_{\mathscr O_{S_0}}
   \mathscr I^{n+1}/\mathscr I^{n+2}\right)=0,
 \]
 \textit{则任意两个这样的提升在 \(X\) 的某个内自同构下共轭,
 此内自同构由一个 \(S\) 上的截面定义,
 该截面在 \(S_0\) 上诱导 \(X_0\) 的单位截面.}
\end{enumerate}

\medskip
\noindent\SGAInlineRefTarget{sga3:iii:lemma:2-7}\textbf{引理 2.7.}
\textit{设 \(S\) 仿射, \(G\) 是仿射 \(S\)-群,
\(\mathscr F\) 是拟凝聚 \(\mathscr O_S\)-模,
\(\mathscr L\) 是局部自由 \(\mathscr O_S\)-模.
假设 \(G\) 按第 I 篇的意义作用于 \(\mathscr F\),
因而也作用于 \(\mathscr F\otimes_{\mathscr O_S}\mathscr L\).
设 \(\Lambda\) 是 \(S\) 的环, \(L\) 是定义
\(\mathscr L\) 的投射 \(\Lambda\)-模. 则有典范同构}
\[
 H^*(G,\mathscr F\otimes_{\mathscr O_S}\mathscr L)
 \simeq H^*(G,\mathscr F)\otimes_\Lambda L.
\]
\footnote{编者注: 这里赋予 \(\mathscr L\) 平凡的 \(G\) 作用.}

设 \(\mathscr A\) 是 \(\mathscr O_S\)-代数
\(\mathscr A(G)\), 并考虑拟凝聚 \(\mathscr O_S\)-模的复形 \(\mathscr C\)
\[
 0\longrightarrow\mathscr F
 \longrightarrow\mathscr F\otimes\mathscr A
 \longrightarrow\mathscr F\otimes\mathscr A\otimes\mathscr A
 \longrightarrow\cdots.
\]
\footnote{编者注: 原证明已扩充并简化. 它还使用了同构
 \(H^n(\mathscr C\otimes\mathscr L)\simeq
 H^n(\mathscr C)\otimes\mathscr L\) 及
\(H_n(\Gamma(-))\simeq\Gamma(H_n(-))\), 其中
\(H^n(\mathscr C)\) 表示复形 \(\mathscr C\) 的上同调层.}

% Source page: Exp3 local p. 36; combined p. 146; printed p. 136.
由 I, \SGARefLink{sga3:i:subsection:5-3}{5.3},
\(H^*(G,\mathscr F)\), 分别地,
\(H^*(G,\mathscr F\otimes\mathscr L)\), 是
\(\Gamma(S,\mathscr C)\), 分别地,
\(\Gamma(S,\mathscr C\otimes\mathscr L)\) 的上同调.
由于 \(S\) 仿射, EGA I, 1.3.12 给出
\[
 \Gamma(S,\mathscr C\otimes\mathscr L)
 \simeq\Gamma(S,\mathscr C)\otimes_\Lambda L.
\]
由于 \(L\) 投射, 因而平坦,
\[
 H^*(\Gamma(S,\mathscr C)\otimes_\Lambda L)
 \simeq H^*(\Gamma(S,\mathscr C))\otimes_\Lambda L,
\]
这就证明了引理.

利用此引理, \SGARefLink{sga3:iii:corollary:2-6}{推论 2.6} 成为:

\medskip
\noindent\SGAInlineRefTarget{sga3:iii:corollary:2-8}\textbf{推论 2.8.}
\textit{设 \(S\) 仿射, \(\mathscr I\) 是定义 \(S_0\) 的幂零理想.
假定每个 \(\mathscr I^{n+1}/\mathscr I^{n+2}\) 在 \(S_0\) 上局部自由.
设 \(Y\) 是在 \(S\) 上平坦的仿射 \(S\)-群,
\(X\) 是平滑 \(S\)-群, \(f_0:Y_0\to X_0\) 是 \(S_0\)-群态射.}
\begin{enumerate}
\item[\(i\)]\SGAInlineRefTarget{sga3:iii:corollary:2-8:item:i}
\textit{若 \(H^2(Y_0,\operatorname{Lie}(X_0/S_0))=0\),
则 \(f_0\) 可提升为 \(S\)-群态射 \(Y\to X\).}
\item[\SGAInlineRefTarget{sga3:iii:corollary:2-8:item:ii}(ii)]
\textit{若 \(H^1(Y_0,\operatorname{Lie}(X_0/S_0))=0\),
则任意两个这样的提升在 \(X\) 的某个内自同构下共轭,
此内自同构由一个 \(S\) 上的截面定义,
该截面在 \(S_0\) 上诱导 \(X_0\) 的单位截面.}
\end{enumerate}

取 \(Y=X\), 得到:

\medskip
\noindent\SGAInlineRefTarget{sga3:iii:corollary:2-9}\textbf{推论 2.9.}
\textit{设 \(S\) 和 \(S_0\) 如上, \(X\) 是平滑仿射 \(S\)-群.}
\begin{enumerate}
\item[\(i\)] \textit{若
\(H^1(X_0,\operatorname{Lie}(X_0/S_0))=0\),
则 \(X\) 在 \(S\) 上的每个在 \(X_0\) 上诱导恒等的自同态,
都是由一个 \(S\) 上的截面定义的内自同构,
该截面在 \(S_0\) 上诱导 \(X_0\) 的单位截面.}
\item[(ii)] \textit{若
\(H^2(X_0,\operatorname{Lie}(X_0/S_0))=0\),
则 \(X_0\) 的每个 \(S_0\)-自同构都可延拓为
\(X\) 的 \(S\)-自同构.}
\footnote{编者注: 设 \(f_0\) 是 \(X_0\) 的
\(S_0\)-自同构, 其逆为 \(g_0\). 由
\SGARefLink{sga3:iii:corollary:2-8:item:i}{2.8(i)}, 将
\(f_0\) 和 \(g_0\) 提升为 \(X\) 的 \(S\)-自同态 \(f\) 和 \(g\).
则 \(g\circ f\) 和 \(f\circ g\) 在 \(X_0\) 上诱导恒等,
所以由 \SGARefLink{sga3:iii:corollary:0-7}{0.7}, 它们是 \(S\)-自同构.
因而 \(f\) 和 \(g\) 也都是 \(S\)-自同构.}
\end{enumerate}

\medskip
\noindent\SGAInlineRefTarget{sga3:iii:remark:2-10}\textbf{注 2.10.}
由 \SGARefLink{sga3:iii:theorem:2-1}{定理 2.1},
涉及 \(H^1\) 的断言具有逆命题.
例如, 设 \(S=I_{S_0}\) 是 \(S_0\) 上的对偶数概形
(II, \SGARefLink{sga3:ii:definition:2-1}{2.1}),
并设 \(X\) 是平坦 \(S\)-群, 使得每个 \(S\) 上的在
\(S_0\) 上诱导恒等的自同构, 都是由一个 \(S\) 上的截面定义的内自同构,
该截面诱导 \(X_0\) 的单位截面. 则
\[
 H^1(X_0,\operatorname{Lie}(X_0/S_0))=0.
\]
\footnote{编者注: 这用于 \SGARefLink{sga3:xxiv:1.13}{第 XXIV 篇, 1.13}.}

\medskip
\noindent\SGAInlineRefTarget{sga3:iii:corollary:2-11}\textbf{推论 2.11.}
\textit{设 \(S,\mathscr I,\mathscr J\) 如
\SGARefLink{sga3:iii:theorem:2-1}{2.1} 中所述.
设 \(Y\) 是在 \(S\) 上平坦的 \(S\)-群概形,
\(X\) 是 \(S\)-群概形, \(f:Y\to X\) 是 \(S\)-群态射.
从 \(f\) 通过与如下元素 \(x\in X(S)\) 共轭而得到的从 \(Y\)
到 \(X\) 的态射所成的集合, 同构于下面的商集;
这些元素诱导 \(X(S_{\mathscr J})\) 的单位元:}
\[
 \mathcal E=
 \operatorname{Hom}_{\mathscr O_{S_0}}(\omega^1_{X_0/S_0},\mathscr J)
 \big/
 \operatorname{Hom}_{\mathscr O_{S_0}}(\omega^1_{X_0/S_0},\mathscr J)
 ^{\operatorname{ad}(Y_0)}.
\]
\textit{这里第二个群由如下 \(\mathscr O_{S_0}\)-态射
\(\omega^1_{X_0/S_0}\to\mathscr J\) 组成: 它们在每个基变换
\(S'\to S_0\) 后给出的态射}
\[
 \omega^1_{X_0/S_0}\otimes_{\mathscr O_{S_0}}\mathscr O_{S'}
 \longrightarrow
 \mathscr J\otimes_{\mathscr O_{S_0}}\mathscr O_{S'}
\]
\textit{在 \(Y_0(S')\) 对第一个因子的作用下不变.}

由 \SGARefLink{sga3:iii:theorem:2-1:item:iii}{定理 2.1(iii)},
此集合为 \(\Gamma(L_0)/H^0(Y_0,L_0)\). 而
\[
 \Gamma(L_0)=
 \operatorname{Hom}_{\mathscr O_{S_0}}(\omega^1_{X_0/S_0},\mathscr J),
\]
且 \(H^0(Y_0,L_0)\) 恰好是刚才所述意义下的
\(\Gamma(L_0)^{\operatorname{ad}(Y_0)}\).

% Source page: Exp3 local p. 37; combined p. 147; printed p. 137.
\medskip
\noindent\SGAInlineRefTarget{sga3:iii:corollary:2-12}\textbf{推论 2.12.}
\textit{在 \SGARefLink{sga3:iii:corollary:2-11}{2.11} 的假设下,
再假设 \(\omega^1_{X_0/S_0}\) 局部自由且秩有限. 则}
\[
 \mathcal E\simeq
 \Gamma\!\left(S_0,
  \operatorname{Lie}(X_0/S_0)\otimes_{\mathscr O_{S_0}}\mathscr J
 \right)
 \big/
 H^0\!\left(Y_0,
  \operatorname{Lie}(X_0/S_0)\otimes_{\mathscr O_{S_0}}\mathscr J
 \right).
\]
事实上,\footnote{编者注: 下面这句话是新增的.}
由局部自由且秩有限得到
\[
 \operatorname{Hom}_{\mathscr O_{S_0}}
 (\omega^1_{X_0/S_0},\mathscr J)
 \simeq
 \operatorname{Lie}(X_0/S_0)\otimes_{\mathscr O_{S_0}}\mathscr J.
\]

\medskip
\noindent\SGAInlineRefTarget{sga3:iii:corollary:2-13}\textbf{推论 2.13.}
\textit{再假设 \(Y_0\) 可对角化. 则}
\[
 \mathcal E\simeq
 \Gamma\!\left(S_0,
  \operatorname{Lie}(X_0/S_0)\otimes_{\mathscr O_{S_0}}\mathscr J
 \right)
 \big/
 \Gamma\!\left(S_0,
  \operatorname{Lie}(X_0/S_0)^{\operatorname{ad}(Y_0)}
  \otimes_{\mathscr O_{S_0}}\mathscr J
 \right),
\]
\textit{其中
\(\operatorname{Lie}(X_0/S_0)^{\operatorname{ad}(Y_0)}\) 是
I, \SGARefLink{sga3:i:proposition:4-7-3}{4.7.3} 的分解中
与 \(Y_0\) 的零特征标对应的因子.}

事实上, 若 \(Y_0\simeq D_{S_0}(M)\), 则所引结果给出直和分解
\[
 \operatorname{Lie}(X_0/S_0)
 =\operatorname{Lie}(X_0/S_0)_0
 \oplus
 \bigoplus_{\substack{m\in M\\m\ne0}}
 \operatorname{Lie}(X_0/S_0)_m.
\]
与 \(\mathscr J\) 作张量积, 得到
\(\operatorname{Lie}(X_0/S_0)\otimes\mathscr J\) 的类似分解, 因而
\[
 H^0\!\left(Y_0,
  \operatorname{Lie}(X_0/S_0)\otimes\mathscr J\right)
 \simeq
 \Gamma\!\left(S_0,
  \operatorname{Lie}(X_0/S_0)_0
  \otimes_{\mathscr O_{S_0}}\mathscr J\right).
\]

\medskip
\noindent\SGAInlineRefTarget{sga3:iii:corollary:2-14}\textbf{推论 2.14.}
\textit{再假设 \(S_0\) 仿射. 则}
\[
 \mathcal E\simeq
 \Gamma\!\left(
  S_0,
  \left[
   \operatorname{Lie}(X_0/S_0)
   \big/
   \operatorname{Lie}(X_0/S_0)^{\operatorname{ad}(Y_0)}
  \right]
  \otimes_{\mathscr O_{S_0}}\mathscr J
 \right).
\]

% Continue with section 3 on Exp3 local p. 37 / combined p. 147.

% Translation coverage: entire component, Section 2 through Corollary 2.14.
% Source lines: 582; translated body lines before this audit comment: 594.
% Math preserved: 396 inline formulas, 30 display blocks, 2 align* blocks, 1 equation; prose \text arguments translated.
% Reference targets: 25/25; reference links: 60/60; equation labels: 2/2; footnotes: 13/13.
% Unresolved TODOs: 0; compile and visual review delegated to the volume integrator.
